% !PRINTED 124
\weberpage{124}
%<<CONT
und man sieht dann, dass diese Elemente alle von den Elementen $C$ verschieden sind, da, wenn $g'^{-1}c'g' = g^{-1}cg$ wäre, auch $c' = g'g^{-1}cgg'^{-1}$ wäre, und es würde also gegen die Voraussetzung $c'$ schon in $C$ vorkommen. Auf diese Weise fährt man mit der Bildung der Gruppen $B, B', B'', \ldots$ und der Systeme $C, C', C'', \ldots$ fort, bis alle Elemente von $P$ in den Systemen $A, C, C', C'', \ldots$ untergebracht sind. Man hat dann die Zahlengleichungen:
\weq{9}{n = \nu j = \mu\varepsilon = \mu'\varepsilon' = \mu''\varepsilon''\ldots}
\weq{10}{n = \nu + \varepsilon + \varepsilon' + \varepsilon'' + \cdots}

Auf Grund dieser Formeln lässt sich nun der \gesperrt{Sylow'sche} Satz durch vollständige Induction beweisen, wobei zwei Fälle zu unterscheiden sind:

1) \gesperrt{Der Grad} $\nu$ \gesperrt{von} $A$ \gesperrt{ist durch} $p$ \gesperrt{theilbar.} In diesem Falle giebt es, weil $A$ eine Abel'sche Gruppe ist, nach \S.\,9, 2. ein Element $a$ in $A$ vom Grade $p$, und die cyklische Gruppe $p^{\mathrm{ten}}$ Grades
\[ 1,\ a,\ a^2,\ \ldots,\ a^{p-1}, \]
die wir mit $Q$ bezeichnen, ist ein Normaltheiler von $P$, weil ja für ein Element $a$ von $A$ immer $x^{-1}ax = a$ sein sollte. Die complementäre Gruppe zu $Q$, $P/Q$ ist vom Grade $n:p$, und ihr Grad ist also kleiner als $n$ und durch $p^{\varkappa-1}$ theilbar. Nach unserer Voraussetzung giebt es also einen Theiler von $P/Q$ vom Grade $p^{\varkappa-1}$, und folglich giebt es nach \S.\,6, 2. einen Theiler von $P$ vom Grade $p^{\varkappa}$.

2) \gesperrt{Der Grad} $\nu$ \gesperrt{von} $A$ \gesperrt{ist nicht durch} $p$ \gesperrt{theilbar.} Da $n$ durch $p$ theilbar ist, so können in diesem Falle nach (10) nicht alle Zahlen $\varepsilon, \varepsilon', \varepsilon'', \ldots$, die alle grösser als 1 sind, durch $p$ theilbar sein. Wenn nun $\varepsilon$ nicht durch $p$ theilbar ist, so ist $\mu$ durch $p^{\varkappa}$ theilbar. Der Grad der Gruppe $B$ ist kleiner als $n$ und durch $p^{\varkappa}$ theilbar, und also giebt es nach Voraussetzung einen Theiler von $B$, der also auch Theiler von $P$ ist, vom Grade $p^{\varkappa}$.

Hiermit ist der Satz I bewiesen.

Zu diesem Satze kommen aber noch wesentliche Ergänzungen.
%CONT>>
